% ============================================================
\chapter{Convex Functions}
\label{ch:convex_functions}
% ============================================================

\section{Introduction}

``The great watershed in optimization is not between linearity and nonlinearity, but between convexity and nonconvexity,'' wrote R.~Tyrrell Rockafellar in 1993. This deceptively simple observation captures a profound truth: convex functions enjoy a property so powerful that it transforms the entire landscape of optimization. Every local minimum is automatically global. There are no hidden valleys, no deceptive plateaus, no traps. If you find a point where the function cannot decrease locally, you have found the best point, period.

The theory of convex functions was forged over more than a century, from the inequalities of Jensen (1906) and the geometric insights of Minkowski, through the duality framework of Fenchel in the 1940s, to the monumental treatise of Rockafellar in 1970. What emerges is a remarkably elegant structure: convex functions can be characterised algebraically (via inequalities), analytically (via derivatives), and geometrically (via epigraphs)---and all three viewpoints converge to the same theory. This chapter builds that theory from the ground up.

\section{Definitions and First Properties}

\begin{definition}[Convex Function]
A function $f : \R^n \to \R \cup \{+\infty\}$ is \emph{convex} if its effective
domain $\mathrm{dom}(f) = \{x \mid f(x) < +\infty\}$ is convex and if for all
$x, y \in \mathrm{dom}(f)$ and all $\theta \in [0,1]$:
\[
    f(\theta x + (1-\theta)y) \leq \theta f(x) + (1-\theta)f(y).
\]
\end{definition}

\begin{definition}[Strict Convexity]
$f$ is \emph{strictly convex} if the above inequality is strict for $x \neq y$ and
$\theta \in (0,1)$.
\end{definition}

\begin{definition}[Strong Convexity]
$f$ is \emph{strongly convex} with parameter $m > 0$ (or $m$-strongly convex) if
$f(x) - \frac{m}{2}\norm{x}^2$ is convex, i.e., for all $x, y$ and $\theta \in [0,1]$:
\[
    f(\theta x + (1-\theta)y) \leq \theta f(x) + (1-\theta)f(y)
    - \frac{m}{2}\theta(1-\theta)\norm{x-y}^2.
\]
\end{definition}

\begin{keyformulas}[title=\textbf{Convexity Hierarchy}]
\[
    \text{strongly convex} \;\Rightarrow\; \text{strictly convex}
    \;\Rightarrow\; \text{convex}.
\]
The reverse implications are false in general.
\end{keyformulas}

\section{Epigraph and Characterization}

\begin{definition}[Epigraph]
The \emph{epigraph} of a function $f : \R^n \to \R \cup \{+\infty\}$ is:
\[
    \mathrm{epi}(f) = \{(x, t) \in \R^n \times \R \mid f(x) \leq t\}.
\]
\end{definition}

\begin{theorem}[Epigraphic Characterization]
\label{thm:epigraph}
$f$ is convex if and only if $\mathrm{epi}(f)$ is a convex subset of $\R^{n+1}$.
\end{theorem}

\begin{proof}
$(\Rightarrow)$ Let $(x_1, t_1), (x_2, t_2) \in \mathrm{epi}(f)$ and
$\theta \in [0,1]$. Then:
\[
    f(\theta x_1 + (1-\theta)x_2) \leq \theta f(x_1) + (1-\theta)f(x_2)
    \leq \theta t_1 + (1-\theta)t_2,
\]
so $(\theta x_1 + (1-\theta)x_2, \theta t_1 + (1-\theta)t_2) \in \mathrm{epi}(f)$.

$(\Leftarrow)$ Let $x_1, x_2 \in \mathrm{dom}(f)$. Then
$(x_1, f(x_1)), (x_2, f(x_2)) \in \mathrm{epi}(f)$. By convexity of the epigraph:
\[
    (\theta x_1 + (1-\theta)x_2, \theta f(x_1) + (1-\theta)f(x_2)) \in \mathrm{epi}(f),
\]
which means $f(\theta x_1 + (1-\theta)x_2) \leq \theta f(x_1) + (1-\theta)f(x_2)$.
\end{proof}

\section{Sublevel Sets}

\begin{definition}[Sublevel Set]
The \emph{$\alpha$-sublevel set} of $f$ is:
\[
    S_\alpha = \{x \in \mathrm{dom}(f) \mid f(x) \leq \alpha\}.
\]
\end{definition}

\begin{proposition}
If $f$ is convex, then all its sublevel sets are convex.
\end{proposition}

\begin{warning}[title=\textbf{Converse is False}]
The converse is false: $f(x) = e^x$ has convex sublevel sets (intervals), but
$g(x) = \sqrt{\abs{x}}$ also has convex sublevel sets without being convex. A function
with convex sublevel sets is called \emph{quasiconvex}.
\end{warning}

\section{Differential Characterizations}

\subsection{First-Order Condition}

\begin{theorem}[First-Order Condition]
\label{thm:first_order}
Let $f : \R^n \to \R$ be differentiable on an open convex set $\Omega$. Then $f$ is
convex on $\Omega$ if and only if:
\[
    f(y) \geq f(x) + \inner{\nabla f(x)}{y - x}, \quad \forall x, y \in \Omega.
\]
\end{theorem}

\begin{proof}
$(\Rightarrow)$ By convexity, for $t \in (0,1]$:
\[
    f(x + t(y-x)) \leq f(x) + t(f(y) - f(x)).
\]
So $\frac{f(x+t(y-x)) - f(x)}{t} \leq f(y) - f(x)$. Taking $t \to 0^+$:
$\inner{\nabla f(x)}{y-x} \leq f(y) - f(x)$.

$(\Leftarrow)$ Let $x, y \in \Omega$ and $z = \theta x + (1-\theta)y$. Then:
\begin{align*}
    f(x) &\geq f(z) + \inner{\nabla f(z)}{x - z}, \\
    f(y) &\geq f(z) + \inner{\nabla f(z)}{y - z}.
\end{align*}
Multiplying by $\theta$ and $1-\theta$ respectively and summing:
\[
    \theta f(x) + (1-\theta)f(y) \geq f(z) + \inner{\nabla f(z)}{\theta x + (1-\theta)y - z} = f(z).
\]
\end{proof}

\begin{intuition}
The first-order condition means that the graph of $f$ always lies above its tangent
hyperplanes. Geometrically, the surface $z = f(x)$ curves upward.
\end{intuition}

\subsection{Second-Order Condition}

\begin{theorem}[Second-Order Condition]
\label{thm:second_order}
Let $f : \R^n \to \R$ be twice differentiable on an open convex set $\Omega$. Then:
\begin{enumerate}
    \item $f$ is convex on $\Omega$ $\Leftrightarrow$ $\nabla^2 f(x) \succeq 0$
          for all $x \in \Omega$.
    \item $f$ is strongly convex with parameter $m$ $\Leftrightarrow$
          $\nabla^2 f(x) \succeq mI$ for all $x \in \Omega$.
    \item If $\nabla^2 f(x) \succ 0$ for all $x$, then $f$ is strictly convex
          (the converse is false).
\end{enumerate}
\end{theorem}

\begin{proof}[Proof of (1), $\Rightarrow$]
By the first-order condition, for all $h \in \R^n$ and small enough $t > 0$:
\[
    f(x + th) \geq f(x) + t\inner{\nabla f(x)}{h}.
\]
By Taylor expansion:
\[
    f(x+th) = f(x) + t\inner{\nabla f(x)}{h} + \frac{t^2}{2}h^T\nabla^2 f(x)h + o(t^2).
\]
Hence $h^T\nabla^2 f(x)h + o(t^2)/t^2 \geq 0$. Letting $t \to 0$:
$h^T\nabla^2 f(x)h \geq 0$.
\end{proof}

\section{Jensen's Inequality}

\begin{theorem}[Jensen's Inequality]
Let $f$ be convex and $x_1, \ldots, x_k \in \mathrm{dom}(f)$ with
$\theta_1, \ldots, \theta_k \geq 0$, $\sum \theta_i = 1$. Then:
\[
    f\left(\sum_{i=1}^k \theta_i x_i\right) \leq \sum_{i=1}^k \theta_i f(x_i).
\]
\end{theorem}

\begin{corollary}[Probabilistic Jensen]
Let $X$ be a random variable taking values in $\mathrm{dom}(f)$ with
$\mathbb{E}[\abs{X}] < \infty$. If $f$ is convex:
\[
    f(\mathbb{E}[X]) \leq \mathbb{E}[f(X)].
\]
\end{corollary}

\section{Operations Preserving Convexity of Functions}

\begin{theorem}[Operations on Convex Functions]
\label{thm:function_operations}
\begin{enumerate}
    \item \textbf{Nonnegative weighted sum}: if $f_1, \ldots, f_k$ are convex and
          $\alpha_1, \ldots, \alpha_k \geq 0$, then $\sum \alpha_i f_i$ is convex.
    \item \textbf{Pointwise maximum}: if $f_1, \ldots, f_k$ are convex, then
          $g(x) = \max_i f_i(x)$ is convex.
    \item \textbf{Supremum}: if $(f_\alpha)_{\alpha \in I}$ is a family of convex
          functions, then $g(x) = \sup_\alpha f_\alpha(x)$ is convex.
    \item \textbf{Affine composition}: if $f$ is convex, then $g(x) = f(Ax + b)$
          is convex.
    \item \textbf{Infimal convolution}: if $f, g$ are convex, then
          $(f \mathop{\square} g)(x) = \inf_y \{f(y) + g(x-y)\}$ is convex.
\end{enumerate}
\end{theorem}

\section{Convexity and Smoothness}

\begin{definition}[$L$-Smoothness]
A differentiable convex function $f$ is \emph{$L$-smooth} if its gradient is
$L$-Lipschitz:
\[
    \norm{\nabla f(x) - \nabla f(y)} \leq L\norm{x - y}, \quad \forall x, y.
\]
\end{definition}

\begin{theorem}[Characterizations of $L$-Smoothness]
\label{thm:smoothness}
For convex and differentiable $f$, the following are equivalent:
\begin{enumerate}
    \item $\nabla f$ is $L$-Lipschitz.
    \item $f(y) \leq f(x) + \inner{\nabla f(x)}{y-x} + \frac{L}{2}\norm{y-x}^2$ for all $x, y$.
    \item $f(y) \geq f(x) + \inner{\nabla f(x)}{y-x} + \frac{1}{2L}\norm{\nabla f(y) - \nabla f(x)}^2$ for all $x, y$.
    \item $\inner{\nabla f(x) - \nabla f(y)}{x - y} \geq \frac{1}{L}\norm{\nabla f(x) - \nabla f(y)}^2$ for all $x, y$.
\end{enumerate}
If $f$ is twice differentiable, these are equivalent to $\nabla^2 f(x) \preceq LI$.
\end{theorem}

\begin{keyformulas}[title=\textbf{Fundamental Sandwich Bound}]
For $f$ convex, $m$-strongly convex, and $L$-smooth:
\[
    f(x) + \inner{\nabla f(x)}{y-x} + \frac{m}{2}\norm{y-x}^2
    \leq f(y) \leq
    f(x) + \inner{\nabla f(x)}{y-x} + \frac{L}{2}\norm{y-x}^2.
\]
The condition number $\kappa = L/m$ measures problem difficulty.
\end{keyformulas}

\section{Important Examples}

\begin{example}[Classical Convex Functions]
\begin{enumerate}
    \item \textbf{Affine functions}: $f(x) = \inner{a}{x} + b$ (convex and concave).
    \item \textbf{Norms}: $\norm{\cdot}_p$ for $p \geq 1$ (convex).
    \item \textbf{Quadratic}: $f(x) = \frac{1}{2}x^TQx + q^Tx + c$ with $Q \succeq 0$.
          Strongly convex iff $Q \succ 0$, with $m = \lambda_{\min}(Q)$, $L = \lambda_{\max}(Q)$.
    \item \textbf{Exponential}: $f(x) = e^{ax}$ (convex on $\R$).
    \item \textbf{Log-sum-exp}: $f(x) = \log(\sum_i e^{x_i})$ (convex on $\R^n$).
    \item \textbf{Negative entropy}: $f(x) = \sum_i x_i \log x_i$ on $\R^n_{++}$.
    \item \textbf{Logistic loss}: $f(x) = \log(1 + e^{-x})$ (convex, smooth).
\end{enumerate}
\end{example}

\section{Continuity of Convex Functions}

\begin{theorem}[Continuity]
\label{thm:continuity}
Every convex function $f : \R^n \to \R$ (finite everywhere) is continuous. More
precisely, every proper convex function is continuous on the interior of its effective
domain.
\end{theorem}

\begin{theorem}[Almost Everywhere Differentiability]
Every convex function $f : \R^n \to \R$ is differentiable almost everywhere (in the
Lebesgue sense). In dimension 1, $f$ has left and right derivatives at every point.
\end{theorem}

\section{Convex Conjugates --- Preview}

\begin{definition}[Fenchel Conjugate]
The \emph{Fenchel conjugate} (or Legendre--Fenchel transform) of
$f : \R^n \to \R \cup \{+\infty\}$ is:
\[
    f^*(y) = \sup_{x \in \R^n} \left\{\inner{y}{x} - f(x)\right\}.
\]
\end{definition}

\begin{remark}
$f^*$ is always convex and lower semicontinuous (as a supremum of affine functions
in $y$), even if $f$ is not. Convex conjugation will be studied in detail in
Chapter~\ref{ch:fenchel}.
\end{remark}

\section{Python Implementation}

\begin{codeblock}[title=\textbf{Convexity verification and visualization}]
\begin{minted}{python}
import numpy as np
import matplotlib.pyplot as plt

def check_convexity_numerical(f, x_range, n_tests=1000):
    """Numerical convexity check in 1D."""
    for _ in range(n_tests):
        x = np.random.uniform(*x_range)
        y = np.random.uniform(*x_range)
        t = np.random.uniform(0, 1)
        lhs = f(t * x + (1 - t) * y)
        rhs = t * f(x) + (1 - t) * f(y)
        if lhs > rhs + 1e-10:
            return False
    return True

# Examples
functions = {
    'x^2': lambda x: x**2,
    'exp(x)': lambda x: np.exp(x),
    '|x|': lambda x: np.abs(x),
    'log(1+exp(x))': lambda x: np.log(1 + np.exp(x)),
}

fig, axes = plt.subplots(2, 2, figsize=(10, 8))
for ax, (name, f) in zip(axes.flat, functions.items()):
    x = np.linspace(-3, 3, 200)
    ax.plot(x, f(x), 'b-', linewidth=2)
    # Tangent at x=1
    x0 = 1.0
    h = 1e-5
    grad = (f(x0 + h) - f(x0 - h)) / (2 * h)
    tangent = f(x0) + grad * (x - x0)
    ax.plot(x, tangent, 'r--', alpha=0.7)
    is_cvx = check_convexity_numerical(f, (-3, 3))
    ax.set_title(f'{name} (convex: {is_cvx})')
    ax.grid(True, alpha=0.3)
plt.tight_layout()
plt.savefig('convex_functions.pdf')
\end{minted}
\end{codeblock}

\section{Exercises}

\begin{exercise}[$\star$]
Show that $f(x) = \max(x_1, \ldots, x_n)$ is convex on $\R^n$.
\end{exercise}

\begin{exercise}[$\star$]
Show that if $f$ is convex and $g$ is affine, then $f \circ g$ is convex.
\end{exercise}

\begin{exercise}[$\star\star$]
Show that a differentiable function $f : \R^n \to \R$ is $m$-strongly convex if and
only if $\inner{\nabla f(x) - \nabla f(y)}{x - y} \geq m\norm{x - y}^2$ for all $x, y$.
\end{exercise}

\begin{exercise}[$\star\star$]
Show that $f(x) = \log\left(\sum_{i=1}^n e^{x_i}\right)$ is convex and compute its
smoothness constant $L$ with respect to the $\ell_2$ norm.
\end{exercise}

\begin{exercise}[$\star\star$]
Let $f$ be convex, proper, and lower semicontinuous. Show that $f$ is bounded below
by an affine function.
\end{exercise}

\begin{exercise}[$\star\star\star$]
Prove the equivalence of the four characterizations of $L$-smoothness in
Theorem~\ref{thm:smoothness}.
\end{exercise}

\begin{exercise}[$\star\star\star$]
Let $f$ be convex on $\R^n$ and $x^*$ a minimizer. Show that if $f$ is $m$-strongly
convex:
\[
    f(x) - f(x^*) \geq \frac{m}{2}\norm{x - x^*}^2, \quad \forall x \in \R^n.
\]
\end{exercise}
